\section*{Chapitre \ref{ch:b2:laser} --- Le laser~: émission stimulée et faisceaux gaussiens}

\begin{solution}{exo:b2:laser:1}
$\nu = \qty{4.74e14}{Hz}$~; $h\nu = \qty{3.14e-19}{J} = \qty{1.96}{eV}$~;
$\qty{3.2e15}{}$ photons par seconde~; $c/2L = \qty{500}{MHz}$~; trois
modes.
\end{solution}

\begin{solution}{exo:b2:laser:2}
$\eu^{-80} \approx 2 \times 10^{-35}$~; égalité seulement quand $T \to \infty$~;
$N_2 > N_1$ exige $\eu^{-h\nu/k_BT} > 1$, c'est-à-dire $T < 0$ formellement. À
\qty{24}{GHz}, $h\nu = \qty{1e-4}{eV}$~: $N_2/N_1 = 0.996$.
\end{solution}

\begin{solution}{exo:b2:laser:3}
$z_{\text{R}} = \qty{1.24}{m}$~; $\theta = \qty{4.0e-4}{rad}$~; $w(\qty{100}{m}) =
\qty{40}{mm}$~; à la Lune \qty{150}{km}. Après $\times 20$~: $w_0 = \qty{10}{mm}$,
$\theta = \qty{2e-5}{rad}$, $z_{\text{R}} = \qty{500}{m}$, \qty{10.2}{mm} à
\qty{100}{m}, \qty{7.7}{km} à la Lune.
\end{solution}

\begin{solution}{exo:b2:laser:4}
$g_{\text{th}} = (0.005 + 0.0101)/0.2 = \qty{0.076}{m^{-1}} < 0.2$~: il y a
effet \omterm{def:b2:laser:cavity}{laser}. Avec $R_2 = 0.90$~: $(0.005 + 0.053)/0.2 = \qty{0.29}{m^{-1}}$~: il n'y
en a pas.
\end{solution}

\begin{solution}{exo:b2:laser:5}
(a) Voir la \cref{prop:b2:laser:einstein}. (b) \qty{600}{nm}~: $h\nu/k_BT
= 80$, rapport $5 \times 10^{34}$~; \qty{1}{cm}~: $4.8 \times 10^{-3}$ ---
l'\omterm{def:b2:laser:processes}{émission stimulée} domine. (c) Aux fréquences micro-ondes,
l'\omterm{def:b2:laser:processes}{émission spontanée} est négligeable et une faible inversion donne déjà du
gain~; la technologie (cavités, \omterm{prop:b2:guided-waves:plates}{guides d'ondes}) existait. (d) $A \propto
\nu^3$, donc $\tau \propto \lambda^3$~: \qty{10}{\micro s}.
\end{solution}

\begin{solution}{exo:b2:laser:6}
(a) $I_{\text{s}} = h\nu/\sigma\tau = \qty{1.0e5}{W/m^2}$. (b) $R =
g_0/\sigma\tau = \qty{3.3e22}{m^{-3}.s^{-1}}$. (c) $I = I_{\text{s}}(5 - 1) =
\qty{4.2e5}{W/m^2}$~; $P_{\text{out}} = 0.01 \times 4.2 \times 10^5 \times
10^{-6} = \qty{4.2}{mW}$. (d) $I = I_{\text{s}}(\sigma R\tau/g_{\text{th}}
- 1) = (h\nu/g_{\text{th}})(R - g_{\text{th}}/\sigma\tau)$~: linéaire en $R$
au-dessus de $R_{\text{th}} = g_{\text{th}}/\sigma\tau$.
\end{solution}

\begin{solution}{exo:b2:laser:7}
(a) $c/2L > \qty{1.5}{GHz}$~: $L < \qty{10}{cm}$. (b) $\delta\nu = \nu\,\delta L/L
= \qty{1.6}{GHz}$~: trois espacements de modes, toute la largeur de gain --- le mode
balaie la courbe de gain et passe la main à son voisin. (c) Un
chauffage tient la longueur à une petite fraction de $\lambda$, en utilisant
l'équilibre de deux modes comme signal d'erreur. (d) $c/\delta\nu \approx
\qty{100}{km}$.
\end{solution}

\begin{solution}{exo:b2:laser:8}
(a) $w_0' = \qty{10}{\micro m}$, $z_{\text{R}} = \qty{0.5}{mm}$, $I = 2P/\pi
w_0'^2 = \qty{6.4e6}{W/m^2}$. (b) $w_0' = \qty{3.4}{\micro m}$, $I =
\qty{5.5e7}{W/m^2}$~; l'image du Soleil~: $10^3 \times (3/0.15)^2 =
\qty{4e5}{W/m^2}$ --- le pointeur est cent fois pire. (c) La puissance de
l'ampoule s'étale sur $4\pi$ et son image sur la rétine est étendue~; toute la
puissance du pointeur tombe sur quelques micromètres, à la longueur d'onde
de sensibilité maximale. (d) $2z_{\text{R}} = \qty{1}{mm}$.
\end{solution}

\begin{solution}{exo:b2:laser:9}
(a) $(N/2)\,h\nu/\tau = \qty{7.6e8}{W/m^3}$. (b) \qty{760}{W} --- le régime continu
est hors de question~: une lampe flash. (c) $10^{22} \times 1.87 \times
10^{-19}/2.3 \times 10^{-4} = \qty{8e6}{W/m^3}$~: \qty{8}{W} par centimètre
cube. (d) Cent fois moins, et en continu~: quelques watts de
lumière de diode suffisent.
\end{solution}

\begin{solution}{exo:b2:laser:10}
(a) $\lambda/\pi w_0$~: \qty{0.17}{rad} (\qty{9.5}{\degree}) et \qty{0.5}{rad}
(\qty{28}{\degree}). (b) La direction étroite diverge le plus~:
l'ellipse en champ lointain est perpendiculaire à l'émetteur. (c) $f\theta$~:
\qty{0.8}{mm} et \qty{2.5}{mm}. (d) Une paire de prismes anamorphoseurs ou une
lentille cylindrique~; ou un couplage dans une fibre monomode.
\end{solution}

\begin{solution}{exo:b2:laser:11}
(a) Intensité $\times 1.04$ toutes les \qty{2}{ns}. (b) $N = P\,(2L/c)/h\nu =
6.4 \times 10^8$ photons~; $\ln(6.4 \times 10^8)/\ln 1.04 \approx 520$
allers-retours, environ \qty{1}{\micro s}. (c) La saturation du gain~; un mode sous le
seuil perd plus par aller-retour qu'il ne gagne et reste au
niveau spontané. (d) Les modes se partagent une seule inversion~; le plus fort
la sature et affame les autres~; les dérives de longueur et de gain changent
le vainqueur --- la sortie saute de mode en mode.
\end{solution}

\begin{solution}{exo:b2:laser:12}
(a) $2.7 \times 10^{17}$ photons~; \qty{1}{GW}. (b) $\theta = \qty{3.4e-7}{rad}$,
\qty{130}{m}~; avec le seeing \qty{1.8}{km}. (c) Fraction $0.1/\pi(1800)^2 =
10^{-8}$~: $2.6 \times 10^9$ photons~; étalement du retour \qty{1.3e-5}{rad}, rayon
\qty{5}{km}~; le télescope capte $(0.5/5000)^2 = 10^{-8}$~: environ $25$
photons, quelques-uns après l'optique. (d) $c \times \qty{100}{ps}/2 = \qty{1.5}{cm}$~;
$\times 1/\sqrt{10^4}$~: \qty{0.15}{mm}.
\end{solution}

\begin{solution}{pb:b2:laser:1}
\textbf{1.} \qty{1.96}{eV}, \qty{4.74e14}{Hz}.

\textbf{2.} $\qty{0.05}{eV} \approx 1.5\,k_BT$~: l'agitation thermique fournit
le petit déficit~; le transfert est quasi résonnant, donc efficace.

\textbf{3.} $\eu^{-20.66/0.0345} = \eu^{-600} \approx 10^{-260}$~: aucun~; seule la
décharge alimente le niveau.

\textbf{4.} Le niveau inférieur se vide dix fois plus vite que le supérieur
ne se désexcite~: il est toujours presque vide --- un schéma à quatre niveaux.

\textbf{5.} $v^* = \sqrt{2k_BT/m} = \qty{580}{m/s}$~; $\nu v^*/c \approx \qty{0.9}{GHz}$~;
la formule exacte donne \qty{1.5}{GHz}.

\textbf{6.} $g_0 = \qty{0.09}{m^{-1}}$~; $\eu^{0.027}$~: $2.7\%$ par passage.

\textbf{7.} $g_{\text{th}} = -\ln(0.999 \times 0.99)/0.6 = \qty{0.018}{m^{-1}}$~;
marge $5$.

\textbf{8.} \qty{500}{MHz}~; trois.

\textbf{9.} $P_{\text{in}} = \qty{1}{mW}/0.01 = \qty{100}{mW}$~; \qty{1.3e5}{W/m^2}.

\textbf{10.} Le gain sature à $g_{\text{th}}$~; l'excédent devient des
photons de sortie (et de l'\omterm{def:b2:laser:processes}{émission spontanée}, de la chaleur).

\textbf{11.} $\qty{1}{mW}/\qty{7.5}{W} = 1.3 \times 10^{-4}$~: l'essentiel de
l'énergie des électrons va ailleurs que vers le métastable de l'hélium, et le
photon emporte \qty{1.96}{eV} des \qty{20.66}{eV} investis.

\textbf{12.} $2 \times 0.072 \times 0.3 = 0.043$~; $N = 0.1 \times 2 \times 10^{-9}/3.14
\times 10^{-19} = 6.4 \times 10^8$~; $20.3/0.043 \approx 470$ allers-retours~;
\qty{0.9}{\micro s}.

\textbf{13.} $\delta\nu = \nu\,\delta L/L = \qty{1.6}{GHz}$~: le mode traverse
toute la courbe de gain --- le \omterm{def:b2:laser:cavity}{laser} saute sur un autre mode à moins que la longueur ne soit
stabilisée.

\textbf{14.} La polarisation $p$ traverse une surface de Brewster sans
réflexion~; la polarisation $s$ perd quelque $15\%$ par surface et
n'atteint jamais le seuil~: sortie polarisée rectilignement sans frais.

\textbf{15.} $z_{\text{R}} = \qty{0.45}{m}$~; $\theta = \qty{6.7e-4}{rad}$~; $w(\qty{10}{m})
= \qty{6.7}{mm}$~: \qty{13}{mm} de diamètre.

\textbf{16.} $2P/\pi w_0^2 = \qty{7}{kW/m^2}$~: sept soleils.

\textbf{17.} $w_0' = \qty{34}{\micro m}$~; \qty{5.6e5}{W/m^2}.

\textbf{18.} $w_0' = \qty{11}{\micro m}$~; \qty{5e6}{W/m^2}~; la rétine supporte
\qty{1}{mW} pendant les \qty{0.25}{s} du réflexe de défense --- classe 2~: visible,
$\le \qty{1}{mW}$, sans danger parce qu'on cligne.

\textbf{19.} \qty{6.7e-5}{rad}~; $w_0 = \qty{3}{mm}$, $z_{\text{R}} = \qty{45}{m}$,
$w(\qty{1}{km}) = \qty{67}{mm}$~: \qty{13}{cm}.

\textbf{20.} $c/\Delta\nu$~: \qty{300}{m}~; \qty{30}{cm}.

\textbf{21.} $2.7 \times 10^{17}$~; \qty{1}{GW}.

\textbf{22.} \qty{1.8}{km}~; \omterm{def:b2:diffraction:huygens}{diffraction} seule~: \qty{130}{m}.

\textbf{23.} $0.1/\pi(1800)^2 = 10^{-8}$~; $2.6 \times 10^9$ photons.

\textbf{24.} $\lambda/d = \qty{1.3e-5}{rad}$, \qty{5}{km}~; $(0.5/5000)^2 = 10^{-8}$~:
$25$ photons, deux ou trois détectés.

\textbf{25.} \qty{1.5}{cm}~; \qty{0.15}{mm}~; le recul se mesure en quelques
semaines.
\end{solution}
